\chead{INTRODUCTION GENERALE}

\setcounter{page}{1}

\chapter*{Introduction générale}
\addstarredchapter{Introduction générale}


Dans un contexte professionnel marqué par une transformation numérique accélérée, la gestion du recrutement s'impose comme un enjeu stratégique incontournable pour toute organisation soucieuse de sa compétitivité. Face à l'accroissement constant du volume des candidatures et à la complexité grandissante des processus de sélection, l'automatisation du recrutement est passée du statut d'avantage concurrentiel à celui de nécessité opérationnelle. C'est dans cette optique que s'inscrit notre projet de fin d'études, à travers la conception et le développement de SmartRec, une plateforme web intelligente entièrement dédiée à l'optimisation du processus de recrutement.
\par 
\textbf C'est dans cette logique que s'articule notre projet de fin d'études, à travers la conception et le développement de SmartRec, une plateforme web intelligente et évolutive. Ce projet a été mené au sein de l'entreprise PwC TAC Tunisie, avec l'objectif de proposer une solution innovante répondant aux besoins actuels des entreprises en matière de gestion et d'automatisation du recrutement.
\par 
\textbf SmartRec s'articule autour de deux profils utilisateurs aux prérogatives bien définies : le recruteur, qui gère les offres d'emploi, administre les comptes utilisateurs et consulte les dossiers de candidature avec leurs scores respectifs, et le candidat, qui parcourt les offres disponibles et soumet sa candidature via un formulaire interactif intelligent structuré en deux étapes distinctes.
\par 
\textbf Techniquement, SmartRec repose sur une architecture combinant Spring Boot, Angular et une base de données sécurisée, enrichie d'un modèle (NLP via SpaCy) assurant le parsing automatique des CVs et le calcul d'un score sur 100, le tout développé selon la méthodologie agile Scrum.
\par 
\textbf Ce rapport retrace l'ensemble des étapes suivies, de l'identification des besoins à la réalisation finale de la plateforme SmartRec, en mettant en avant les choix techniques, les défis rencontrés et les compétences mobilisées. Il est structuré pour refléter l'évolution du projet : le premier chapitre présente le cadre général du projet, suivi de l'analyse des besoins. Les chapitres suivants décrivent les sprints de développement, chacun portant sur un axe fonctionnel spécifique. Nous débuterons ainsi par une présentation du cadre général du projet, afin d'ancrer le lecteur dans son contexte de réalisation et ses objectifs initiaux.